\par\Needspace{7\baselineskip}\bigskip
\begin{center}
\pdfbookmark[2]{Gauss to Olbers, Göttingen, 15 June 1814}{pallas-brief-18140615}
\textsc{Gauss} to \textsc{Olbers}, Göttingen, 15 June 1814.
\end{center}
\noindent\hspace*{\parindent}. . . . .\quad Things are going strangely with Pallas. \textsc{Nicolai} has finished the numerical calculation of the perturbations by Saturn for the 9 oppositions, but it makes the previous fine agreement somewhat worse. The need to increase Jupiter's mass nevertheless remains certain, and the increase becomes a little greater, namely \(=\frac{1}{40}\mathpunct{\raisebox{.3ex}{.}}\) Without correcting Jupiter's mass, the best possible agreement in longitude is as follows:
\begin{center}
\begingroup\small\setlength{\arraycolsep}{3pt}
\(\underbrace{\begin{array}{@{}rrrrrrrrr@{}}
-6''{,}7&-27''{,}7&-4''{,}2&+28''{,}7&+31''{,}0&-18''{,}7&-6''{,}5&-33''{,}3&+0''{,}1
\end{array}}\)\par
Sum of the squares . . . . . . 4110\par\smallskip
Mean error \(=\BareRoot{\frac{4110}{9}}=21''\).
\endgroup
\end{center}
\par With Jupiter's mass increased by \(\frac{1}{40}\), however:
\begin{center}
\begingroup\small\setlength{\arraycolsep}{3pt}
\(\underbrace{\begin{array}{@{}rrrrrrrrr@{}}
-2''{,}3&-9''{,}8&+11''{,}3&+8''{,}5&+5''{,}1&-8''{,}8&-5''{,}7&-5''{,}7&+7''{,}4
\end{array}}\)\par
Sum of the squares \(=525\), which by a slight refinement can be reduced further to 503, so that the mean error becomes \(=7''{,}5\).
\endgroup
\end{center}
\par The latitudes, on the other hand, agree very well and still a little better than before the perturbations by {\AstroSymbols ♄} were included, as follows:
\begin{center}
\begingroup\small\setlength{\tabcolsep}{3pt}
\begin{tabular}{@{}rrrrrrrrr@{}}
\(+5''{,}8\)&\(-0''{,}1\)&\(-2''{,}0\)&\(-3''{,}2\)&\(-8''{,}5\)\footnotemark&\(-1''{,}7\)&\(-8''{,}2\)&\(+1''{,}2\)&\(-4''{,}3\).
\end{tabular}
\endgroup
\end{center}
\footnotetext{The {\AstroSymbols ☍} poorly observed in latitude.}
\par\SourceLeaf{431}I leave open whether an error may have crept into one of the elements in Mr \textsc{Nicolai}'s calculation of the longitudes. It is certainly strange that if I reverse the sign of all the perturbations of the epoch, the agreement becomes \textls[180]{far} better and the sum of the squares of the errors is only about 200. I have suggested a way for Mr \textsc{Nicolai} to check this. I shall hardly finish the perturbations by Mars in less than \(\frac{1}{2}\) year, although they are very small. I am bent on completing them in such a way that nothing is left to be desired.
\par
